\documentclass[twocolumn]{aastex631}
\begin{document}
\title{The reionization ionizing-photon-budget shortfall for star-forming galaxies is not robust to systematics at z\~{}5 (o32 systematic corner)}
\author{NebulaMind Lab (autonomous pipeline)}
\affiliation{NebulaMind Lab, \url{https://lab.nebulamind.net}}
\begin{abstract}
Reionization ionizing-photon-budget reconciliation at z\~{}5: star-forming galaxies require f\_esc=0.019 (+0.020/-0.010) to close the budget (Madau-Dickinson SFRD, log xi\_ion=25.5+/-0.15, clumping C=2-5, JWST-SFRD tail), versus indirect-proxy-inferred f\_esc=0.080 (+0.146/-0.051) from LzLCS O32/beta calibrations. Median delta(required-inferred)=-0.057 dex-frac (16-84\%: -0.202 to -0.005); 13\% of the systematic MC shows a shortfall. Conclusion: the budget CLOSES within the systematic, and the sign holds under both O32 and beta calibrations. The result is bounded by the xi\_ion x clumping x proxy-calibration systematic, not by statistics. Generated autonomously from public data (jwst) via the NebulaMind Lab runner: a bounded, reproducible, descriptive study.
\end{abstract}
\section{Introduction}
Introduction: Recent studies have highlighted potential discrepancies in the reionization photon budget, suggesting that star-forming galaxies may not produce enough ionizing photons to account for the observed reionization [Muoz2024]. This raises questions about our understanding of the cosmic ionizing photon budget and the role of star-forming galaxies in driving reionization. To address this issue, we revisit the ionizing photon budget calculation using a literature-anchored approach, while acknowledging the limitations and uncertainties associated with this method.
\section{Data and method}
Data and method: We adopt the Madau \& Dickinson (2014) analytic fitting function for the cosmic star formation rate density (SFRD). For the ionizing photon production efficiency (xi\_ion), we use published values with an uncertainty of log xi\_ion = 25.5 ± 0.15, recognizing that these values may vary depending on the measurement techniques and assumptions used in different studies. The escape fraction (f\_esc) is inferred from O32/beta calibrations, specifically those derived from LzLCS [Chisholm+22, Flury+22; Simmonds+24], with the caveat that these calibrations may not fully capture the complexities of ionizing photon escape in star-forming galaxies.
\section{Result}
Result: Our calculations show that star-forming galaxies require an escape fraction of f\_esc = 0.019 (+0.020/-0.010) to close the reionization photon budget at z\~{}5, assuming a clumping factor C between 2 and 5. In contrast, indirect proxy-inferred values from LzLCS O32/beta calibrations suggest an escape fraction of f\_esc = 0.080 (+0.146/-0.051). The median difference between the required and inferred escape fractions is -0.057 dex-frac (16-84\%: -0.202 to -0.005), with 13\% of systematic Monte Carlo simulations showing a shortfall in the photon budget. However, we note that these results are subject to uncertainties due to the limited exploration of the clumping factor range and potential systematic errors introduced by using O32/beta calibrations as proxies for f\_esc.
\begin{figure}\centering\includegraphics[width=\columnwidth]{result.png}\caption{The reionization ionizing-photon-budget shortfall for star-forming galaxies is not robust to systematics at z\~{}5 (o32 systematic corner)}\end{figure}
\section{Caveats}
Caveats: Our analysis relies on literature-anchored values for key parameters, which may introduce uncertainties due to variations in measurement techniques and assumptions across different studies. Additionally, our calculation assumes a specific clumping factor range (2-5), which can significantly impact the required escape fraction. A more comprehensive exploration of the clumping factor's effect on the ionizing photon budget is needed to fully understand its implications. Furthermore, the use of O32/beta calibrations as proxies for f\_esc introduces potential systematic errors, and future work should investigate alternative methods for estimating f\_esc. Finally, our analysis does not account for potential contributions from other sources of ionizing photons, such as active galactic nuclei or X-ray binaries, which could affect the overall photon budget. A more complete understanding of these additional sources is necessary to provide a comprehensive picture of the reionization process.
\section*{References}
\footnotesize
[Muoz2024] Muñoz et al. (2024). Reionization after JWST: a photon budget crisis?. bibcode 2024MNRAS.535L..37M \par [Park2022] Park et al. (2022). Calibrating excursion set reionization models to approximately conserve ionizing photons. bibcode 2022MNRAS.517..192P \par [Duncan2015] Duncan et al. (2015). Powering reionization: assessing the galaxy ionizing photon budget at z < 10. bibcode 2015MNRAS.451.2030D \par [Davies2021] Davies et al. (2021). The Predicament of Absorption-dominated Reionization: Increased Demands on Ionizing Sources. bibcode 2021ApJ...918L..35D \par [Madau2017] Madau (2017). Cosmic Reionization after Planck and before JWST: An Analytic Approach. bibcode 2017ApJ...851...50M

\end{document}
